
\subsubsection*{Repeat the period with one explicit company book}
The next assessment may have several earlier liabilities to preserve. A
procedure that always separates one old liability from one current liability
cannot represent that situation. The common period book instead retains an
ordered collection of stable liability records, an immutable financial and
loss opening, current financial state, assessment histories and archived
periods. Each new period uses the same operations. The checked predecessor is
imported once, and a finite sequence of operations supplies the later events
and explicit assessment information.

A first current assessment adds one liability; a replacement changes that
liability's cumulative assessment while keeping its identity and payments.
Every earlier liability stays separate. Starting after the payment and
depreciation example, cash is $709/2$, the two liabilities are $11/4$ and 3,
and profit is $-2$. A hypothetical assessment of 4 adds a third liability,
raises total payable to $39/4$ and reduces profit to $-6$, without moving cash.
Paying 4 with allocations of 1, 1 and 2 then leaves cash at $701/2$ and the
liabilities at $7/4$, 2 and 2. If a fresh cumulative calculation replaces
the current assessment of 4 with 3, only the difference of $-1$ enters current
tax expense. The third liability becomes 1, total payable becomes $19/4$
and profit becomes $-5$. Replacing the assessment with less than the 2 already
paid requires a receivable or refund method; the selected no-credit model
refuses that operation.

Loss replacement starts from the immutable loss opening for the period. For
example, using 4 from an opening quantity of 9 leaves 5; replacing that
calculation with a use of 3 leaves 6, not 2. The earlier use remains in the
assessment history. Expiry is checked at the stated usage date, and a loss
identity cannot be reused after its earlier generation has been superseded
or carried into another period. Lease-right populations likewise include
identities observed in intermediate states, so a right that later disappears
is not silently removed from the period's tax-information requirements.

Fresh deferred measurement follows the current assessment and binds the actual
financial state. Its item bases, recorded balances, support and shared recovery
capacity must reconcile under the selected model. Another payment, native
event or current assessment makes that state-specific measurement stale. The
earlier calculation remains available, but the next boundary requires a new
measurement of the changed book. This dependency prevents an old recognition
capacity from renewing itself merely because the program continues.

At the financial endpoint, the current assessment and loss cursor must reach
the same exclusive date, the closing book must be dated the preceding day,
and the deferred measurement must bind that exact closing state. The
provisional carry closes actual profit into retained earnings and archives the
period's balances, reports, liabilities, payments, assessments and findings.
In the worked sequence, closing profit of $-5$ changes the debit-positive
retained balance from $53/4$ to $73/4$. Cash remains $701/2$. The next period
opens the three liabilities at $7/4$, 2 and 1, with no new settlements or cash
movements. Both tax assessments begin unassessed.

A new assessment of 2 in that period adds a fourth liability and raises total
payable to $27/4$. Profit is $-2$ and operating cash remains zero. The same
deferred-measurement and carry operations can then close that period, changing
retained earnings to $81/4$ while preserving all four liability identities.
These exact recurrences check the selected reference book and its refusal
behavior. Actual reporting completeness, admission of source evidence and
broader company or sector models still require their own specified methods
and evidence.
